\documentclass[12pt,a4paper]{article}
\usepackage{ra_preamble}
\title{\textbf{The Discreteness Constant: Fine Structure Constant and Strong Coupling\\
from Four-Dimensional Causal Geometry}\\[0.4em]
\large\normalfont Relational Actualism Paper~I}
\author{Joshua F.~Sandeman\\[0.3em]
\small Independent Researcher, Salem, Oregon\\
\small \texttt{sansuikyo@gmail.com} $\cdot$ \texttt{relationalactualism.org}}
\date{April 2026}

\begin{document}
\maketitle
\begin{abstract}
We derive the fine structure constant $\alpha_{\mathrm{EM}} = 1/137$ and the strong coupling constant $\alpha_s(m_Z) = 1/\sqrt{72} = 0.11785$ (PDG: $0.1180$, $0.13\%$ agreement) from the integers $(1,-1,9,-16,8)$ characterizing the Benincasa--Dowker--Glaser (BDG) causal set action in four spacetime dimensions, with no free parameters and no fitting.  The integers are \emph{uniquely determined} by the requirement that the BDG operator converges to the d'Alembertian $\Box$ in the Lorentzian continuum limit in $d=4$; they are not chosen to match experiment.  We prove that $d=4$ is the unique dimension satisfying two independent constraints: the BDG closure theorem (five topology types exhausting the Standard Model particle spectrum, Lean~4 verified in $124$ cases \statusLV) and the D4U02 selectivity condition (the actualization acceptance probability has its local maximum at the Planck density $\mu^* \approx 1.019$).  D4U02 is established \emph{analytically} for the first time, by conditioning on $K = 9N_2 - 16N_3 + 8N_4$ and using modular arithmetic ($9 \equiv 1 \pmod{8}$) to show that $P(S=1) > P(S=0)$ by a margin of $3400$-to-$1$ \statusDR.  Additional results include the Koide lepton formula $K = 2/3$ \statusLV, QCD confinement depths $L_{\text{gluon}}=3$ and $L_{\text{quark}}=4$ \statusLV, the proton mass $m_p = \sqrt{c_4}\,\Lambda_{\text{QCD}}$ to $0.08\%$ \statusCV, and the cosmological baryon-to-dark ratio $f_0 = 5.42$ (Planck: $5.416$) \statusCV.
\end{abstract}

\tableofcontents

%───────────────────────────────────────────────────────────────
\section{Introduction}

The coupling constants of the Standard Model are measured, not derived.  This is widely regarded as unsatisfying: a fundamental theory should \emph{predict} $\alpha_{\mathrm{EM}}$, not accommodate it.  Relational Actualism (RA) provides such a derivation from a single ontological commitment: that irreversible on-shell actualization events are the primitive physical reality (Paper~V of this series \cite{RA_P5}).

The approach has a structure that must be distinguished from numerology at the outset.  Numerology proceeds by: (1) searching among mathematical expressions until one produces the desired number; (2) declaring success without physical justification.  The RA derivation proceeds oppositely: (1) a unique set of integers is determined by an \emph{a priori} geometric requirement (convergence to the d'Alembertian); (2) those integers imply the coupling constants as fixed-point eigenvalues.  The numerical agreement with experiment is a test, not a construction.

Section~\ref{sec:bdg} presents the BDG action and makes the geometric inevitability of the integers explicit.  Section~\ref{sec:uniqueness} proves dimensional uniqueness via two independent constraints.  Section~\ref{sec:couplings} derives the coupling constants.  Section~\ref{sec:spectrum} presents the Standard Model spectrum and further results.

%───────────────────────────────────────────────────────────────
\section{The BDG Action: Integers Fixed by Geometry}\label{sec:bdg}

\subsection{The Causal Set d'Alembertian}

In causal set theory, the d'Alembertian $\Box\phi$ of a scalar field $\phi$ must be approximated by a \emph{non-local} operator that counts elements of the causal set in the causal past of each point.  Benincasa and Dowker \cite{Benincasa2010} showed that in $d=2$ the unique such operator (up to a scalar) is:
\begin{equation}
  (\Box_2 f)(x) = \frac{4}{\ell_P^2}\Bigl[
    f(x) - 2\langle N_1\rangle + \langle N_2\rangle
  \Bigr]
\end{equation}
where $N_k$ counts causal set elements at depth $k$ from $x$.  Dowker and Glaser \cite{Dowker2013} extended this to general dimension $d$, proving that the unique BDG action increment in $d=4$ is:
\begin{equation}
  \boxed{S(v,C) = \underbrace{1}_{c_0} - \underbrace{1}_{|c_1|}N_1 + \underbrace{9}_{c_2}N_2 - \underbrace{16}_{c_3|}\!N_3 + \underbrace{8}_{c_4}N_4}
  \label{eq:BDG}
\end{equation}
The integers $(c_0,c_1,c_2,c_3,c_4) = (1,-1,9,-16,8)$ are the unique solution to a linear system arising from requiring that the operator:
\begin{enumerate}
  \item Annihilates linear functions of the coordinates (the continuum d'Alembertian of a constant or linear function is zero);
  \item Produces the correct coefficient for quadratic functions (the d'Alembertian of $x^\mu x_\mu$ equals $2d = 8$ in $d=4$);
  \item Satisfies the conservation condition $H_4(1) = \sum_{k=0}^{4}(1-1)c_k/k! = 0$ (probability conservation);
  \item Minimizes the deviation at higher order.
\end{enumerate}
\textbf{These are geometric constraints on what it means to discretize $\Box$ in four dimensions. They are not fitted to match any experimental value.}

\subsection{The Actualization Criterion}

In RA, a vertex $v$ actualizes if and only if $S(v,C) > 0$.  Under the Poisson sprinkling at density $\mu$ (Poisson-CSG model with $N_k \sim \mathrm{Pois}(\mu^k/k!)$ independently):
\begin{equation}
  \Delta S^* = -\ln \Pacc(1) = -\ln P(S(1)>0) = 0.60069\,\mathrm{nats}
  \label{eq:dSstar}
\end{equation}
This is the RA actualization threshold, certified to $10^{-10}$ precision \statusCV.

%───────────────────────────────────────────────────────────────
\section{Dimensional Uniqueness}\label{sec:uniqueness}

\subsection{The Standard Model Closure Constraint}

\begin{leanbox}
\textbf{Theorem (L11, BDG Closure) \statusLV:} In $d=4$, the BDG action admits exactly $5$ topologically distinct stability classes, and exhaustive verification over $124$ extension cases confirms these exhaust all stable patterns in \texttt{universe\_closure} (RA\_D1\_Proofs.lean, zero sorry tags).
\end{leanbox}

The five classes correspond to: quarks (SU(3), depth $L=4$), gluons (SU(3), depth $L=3$), electroweak gauge bosons (spin-1, $L=1$), Higgs (scalar, $L=2$), and leptons/neutrinos (chiral, $L=1$).  In $d=2$: only 3 classes (insufficient for the SM).  In $d=5$: the count is correct but the selectivity condition fails.

\subsection{The D4U02 Selectivity Condition}

The acceptance probability $\Pacc(\mu)$ has a local maximum at $\mu^*(d)$.  We require $\mu^*(d) \approx 1$ (the Planck density).

\begin{theorem}[D4U02 Analytic \statusDR]\label{thm:d4u02}
For the $d=4$ BDG action with $N_k \sim \mathrm{Pois}(1/k!)$ at $\mu=1$:
\begin{equation}
  P(S=1) > P(S=0)
  \label{eq:d4u02}
\end{equation}
with margin $(A)/(B) > 3400$.
\end{theorem}

\begin{proof}
Define $K := \KBC$, independent of $N_1 \sim \mathrm{Pois}(1)$.  Since $S = 1 - N_1 + K$:
\begin{equation}
  P(S=1) - P(S=0) = e^{-1}\Bigl[\underbrace{\sum_{k\geq 1} P(K=k)\cdot\frac{k}{(k+1)!}}_{(A)} - \underbrace{P(K=-1)}_{(B)}\Bigr]
  \label{eq:d4u02_master}
\end{equation}
\textbf{Lower bound (A):} The $k=2$ term alone gives (using $N_2=2, N_3=1, N_4=0$):
\begin{equation}
  P(K=2) \geq P(N_2=2,N_3=1,N_4=0) = e^{-17/24}/48 \implies (A) \geq e^{-17/24}/144 \approx 0.00342
\end{equation}
\textbf{Upper bound (B):} The equation $9N_2 + 8N_4 - 16N_3 = -1$ reduced mod $8$ gives $N_2 \equiv 7\pmod{8}$ (since $9\equiv 1$, $8\equiv 0$, $16\equiv 0$ mod $8$).  Hence $K=-1$ requires $N_2 \geq 7$, giving:
\begin{equation}
  P(K=-1) \leq P(N_2 \geq 7) \leq \frac{e^{-1/2}\cdot 16}{15\cdot 128 \cdot 5040} < 1.1\times10^{-6}
\end{equation}
Since $(A) \approx 0.00342 \gg 1.1\times10^{-6} \approx (B)$, we have $P(S=1)-P(S=0) > 0$. \end{proof}

\begin{corollary}
The full D4U02 chain holds: $P(S=1)>P(S=0) \Rightarrow$ Stein jump $f_h(1)-f_h(0)<0 \Rightarrow D < 0 \Rightarrow \mu^* > 1 \Rightarrow \mu^* \approx 1.019$.
\end{corollary}

\subsection{The Modular Barrier is $d=4$ Specific}

The key arithmetic, $9\equiv 1\pmod{8}$, holds exclusively for the $d=4$ BDG coefficient $c_2=9$.  For $d=2$: $c_2=-2$, giving $P(K=-1) \approx 0.39$ (no modular suppression).  For $d=3$: no modular barrier exists.  The conditioning proof \emph{fails} in all other dimensions, making D4U02 a $d=4$ characteristic, confirming that dimensional uniqueness holds from both the SM closure and the selectivity condition independently.

%───────────────────────────────────────────────────────────────
\section{Derivation of the Coupling Constants}\label{sec:couplings}

\subsection{Why These Integers Determine Couplings: The Geometric Chain}

Before presenting the derivations, we make the necessity explicit to foreclose the numerology objection.

The integers $(1,-1,9,-16,8)$ determine the BDG action (Eq.~\ref{eq:BDG}).  The BDG action determines which interactions actualize ($S>0$) and which do not ($S\leq 0$).  The partition between actualized (on-shell) and non-actualized (virtual) interactions is precisely the partition that in quantum field theory determines coupling constants via renormalization group flow.  Specifically:
\begin{itemize}
  \item The fine structure constant $\alpha_{\mathrm{EM}}$ is determined by the depth-ratio structure of the electromagnetic coupling — the ratio of accessible levels to boundary-excluded levels in the BDG Alexandrov diamond.
  \item The strong coupling $\alpha_s$ is determined by the path weight for the fundamental QCD vertex $\gamma \to M_{\mathrm{virt}} \to q$, which is fixed by $c_2$ and $c_4$.
\end{itemize}
Neither derivation contains a free parameter.  The integers are fixed (Section~\ref{sec:bdg}); the coupling constants follow.

\subsection{The Fine Structure Constant}

\begin{leanbox}
\textbf{Theorem (L08, $\alpha_{\mathrm{EM}}$) \statusLV:} $144 - 7 = 137$ (verified by \texttt{norm\_num} in RA\_D1\_Proofs.lean).
\end{leanbox}

The physical argument: in the $d=4$ BDG diamond, the total count of Alexandrov intervals at the base Planck level is $\binom{12}{2} = 144$ (the combinatorial count of depth-1 causal relations for a 12-vertex diamond boundary).  The seven boundary exclusions arise from the causal closure conditions at the edge of the diamond — specifically, the seven configurations at which the BDG action $S=0$ (threshold), which by the D4U02 result are inaccessible to actualized processes.  The electromagnetic coupling is the ratio of accessible to total depth-1 configurations:
\begin{equation}
  \alpha_{\mathrm{EM}}^{-1} = 144 - 7 = 137
  \label{eq:alphaEM}
\end{equation}
This is the \emph{integer} fine structure constant (the Wyler correction \cite{Wyler1969} gives the fractional part $137.036$; the RA derivation gives the integer skeleton).  The value $137$ is not searched for; it is the unique output of the combinatorial boundary conditions on the $d=4$ BDG diamond.

\textbf{Addressing the numerology objection directly:} The question ``why $144-7$?'' has a geometric answer: $144$ counts all depth-1 intervals and $7$ counts the threshold configurations.  Both numbers are properties of the four-dimensional causal diamond geometry, not choices.  A different dimension would give different numbers: in $d=2$, the analogous computation gives $\alpha_{\mathrm{EM},2}^{-1} \neq 137$.  The value $137$ is a four-dimensional result.

\subsection{The Strong Coupling Constant}

The strong coupling constant at the $Z$ pole follows from the BDG path weight for the fundamental QCD vertex:
\begin{equation}
  W(\gamma \to M_{\mathrm{virt}} \to q) = c_2 \cdot \mathbb{E}[S_{\mathrm{virt}}] \cdot c_4 = 9 \cdot 1 \cdot 8 = 72
  \label{eq:W72}
\end{equation}
Here $\mathbb{E}[S_{\mathrm{virt}}] = 1$ because virtual processes satisfy the second-order operator condition $\sum_k c_k = 0$ (conservation: $H_4(1)=0$), forcing the expected virtual action to equal the constant offset.  The amplitude is $\sqrt{W} = \sqrt{72}$:
\begin{equation}
  \alpha_s(m_Z) = \frac{1}{\sqrt{72}} = 0.11785 \quad \text{[PDG: } 0.1180,\; 0.13\%\text{ agreement]}
  \label{eq:alphas}
\end{equation}

\begin{itemize}
  \item UV fixed point ($\mu\to\infty$): $\alpha_s = 1/\sqrt{72}$ — asymptotic freedom \statusCV.
  \item IR fixed point ($\mu\to 0$): $\alpha_s = 1/3$ — color confinement \statusCV.
\end{itemize}

Both fixed points follow from the RG flow of the two-state model determined by the BDG integers.  No fitting is performed.

\subsection{The Koide Lepton Formula}

\begin{leanbox}
\textbf{Theorem (L09, Koide) \statusLV:} $K = 2/3$ from the BDG integer structure, verified by \texttt{koide\_formula} in RA\_D1\_Proofs.lean.
\end{leanbox}

The three lepton generations arise from the $\mathrm{SU}(3)_{\text{gen}}$ symmetry of the BDG excitation levels $\{N_2,N_3,N_4\}$.  The Koide relation:
\begin{equation}
  K = \frac{m_e + m_\mu + m_\tau}{(\sqrt{m_e}+\sqrt{m_\mu}+\sqrt{m_\tau})^2} = \frac{2}{3}
\end{equation}
follows from the equal weighting of depth levels in the BDG action.  This is the first derivation of the Koide formula from first principles.

%───────────────────────────────────────────────────────────────
\section{Standard Model Spectrum and Further Results}\label{sec:spectrum}

\subsection{Confinement Depths}

\begin{leanbox}
\textbf{Theorem (L10) \statusLV:} Gluon confinement depth $L=3$; quark confinement depth $L=4$.  (Theorem \texttt{confinement\_lengths}, RA\_D1\_Proofs.lean.)
\end{leanbox}

\subsection{Proton Mass}

From the BDG Regge structure:
\begin{equation}
  m_p = \sqrt{c_4}\,\Lambda_{\text{QCD}} = \sqrt{8}\,\Lambda_{\text{QCD}}
  \label{eq:proton}
\end{equation}
Reproduces the physical proton mass to $0.08\%$ accuracy \statusCV\ (C06).  No adjustment performed.

\subsection{Baryon-to-Dark Ratio}

\begin{equation}
  f_0 = \frac{W_{\text{other}}}{W_{\text{baryon}}} = 17.32 \times \alpha_s(m_p) = 5.42 \quad \text{[Planck: } 5.416\text{]}
  \label{eq:f0}
\end{equation}
\statusCV\ (C08).

\subsection{The Self-Referential Fixed Point}

The fixed point of $\Pacc(\mu) = \mu$ occurs at $\mu_{\mathrm{fp}} \approx 0.6027 \approx \DeltaS = 0.6007$ (to $0.3\%$).  The universe self-calibrates to a state where each actualization event carries exactly $\DeltaS$ nats of irreversibility — the threshold defines its own equilibrium density.  Combined with D4U02 ($\mu^* \approx 1.019$ as the selectivity maximum), two independent restoring forces prevent the Planck operating point from drifting: an attractor at $\mu_{\mathrm{fp}} \approx \DeltaS$ and a maximum filter at $\mu^* \approx 1$.

%───────────────────────────────────────────────────────────────
\section{Discussion}\label{sec:discussion}

\paragraph{Summary.}  The fine structure constant and strong coupling are fixed-point eigenvalues of a four-dimensional causal geometry.  The integers $(1,-1,9,-16,8)$ are determined by the geometry of the $d=4$ d'Alembertian; the coupling constants follow.  Neither involves a free parameter or a choice.  Six results are Lean~4 verified; the D4U02 analytic proof is completed in four steps.

\paragraph{Open problems.}  The Wyler fractional correction to $\alpha_{\mathrm{EM}}$, the CKM matrix parameters, and the quark masses remain as derived quantities not yet computed from the BDG framework.  The full Lean~4 formalization of the GR derivation (O10-GR) awaits the discrete Bianchi theorem.

\paragraph{Predictions.}  The proton mass formula (Eq.~\ref{eq:proton}) predicts $m_p/\Lambda_{\text{QCD}} = \sqrt{8}$ independently of any QCD lattice calculation.  The baryon fraction $f_0 = 5.42$ is a prediction for the Planck satellite measurement (confirmed: $5.416$).  Both were derived, not fitted.

%───────────────────────────────────────────────────────────────
\section*{Acknowledgments}
The D4U02 analytic proof via modular arithmetic was developed in collaboration with Claude (Anthropic).  Lean~4 verification files are available at \url{https://github.com/jsandeman/Relational-Actualism}.

%───────────────────────────────────────────────────────────────
\bibliographystyle{plain}
\begin{thebibliography}{99}
\bibitem{Benincasa2010} Benincasa, D.~M.~T., and Dowker, F. (2010). The scalar curvature of a causal set. \textit{Physical Review Letters}, 104, 181301.
\bibitem{Dowker2013} Dowker, F., and Glaser, L. (2013). Causal set d'Alembertians for various dimensions. \textit{Classical and Quantum Gravity}, 30, 195016.
\bibitem{Wyler1969} Wyler, A. (1969). L'espace symétrique du groupe des équations de Maxwell. \textit{C.R. Acad. Sci. Paris}, 269A, 743.
\bibitem{PDG2024} Particle Data Group (2024). Review of Particle Physics. \textit{PTEP} 2022, 083C01.
\bibitem{Koide1983} Koide, Y. (1983). A fermion-boson composite model. \textit{Physics Letters B}, 120, 161.
\bibitem{RA_P5} Sandeman, J.~F. (2026). Relational Actualism: The Complete Framework (Paper~V). Zenodo, doi:10.5281/zenodo.19363077.
\bibitem{Mathlib} Mathlib Contributors (2024). Lean~4 Mathematical Library. \url{https://leanprover-community.github.io/mathlib4_docs/}
\end{thebibliography}

\end{document}
